% TOP_PROBLEM: {"id": 2605, "title": "Kourovka Notebook Problem 21.96", "category_id": 17, "category": {"id": 17, "name": "group_theory", "display_name": "Group Theory", "description": "Problems about groups, group actions, representations, and related algebraic structures.", "slug": "group-theory", "order_index": 17, "created_at": "2026-07-31T15:26:25.671Z"}, "status": "open", "problem_number": "KOU-21.96", "set_id": 10}
% FIRST_POSTED: 2026-10-01
% ENTRY_KIND: statement_only
% UnsolvedMath ID 2605; source catalogue code KOU-21.96
% !TeX program = lualatex
% Definition and sources only; this file is not an LLM solution attempt.
\documentclass[11pt,a4paper]{article}
\usepackage[margin=25mm]{geometry}
\usepackage{fontspec}
\setmainfont{lmroman10-regular.otf}[BoldFont=lmroman10-bold.otf,
 ItalicFont=lmroman10-italic.otf,BoldItalicFont=lmroman10-bolditalic.otf]
\usepackage{amsmath,amssymb,mathtools}
\usepackage{unicode-math}
\setmathfont{latinmodern-math.otf}
\AtBeginDocument{\let\setminus\smallsetminus}
\usepackage{xurl,hyperref}
\hypersetup{colorlinks=true,urlcolor=blue}
\setcounter{secnumdepth}{0}
\setlength{\parindent}{0pt}
\setlength{\parskip}{5pt}
\setlength{\emergencystretch}{3em}
\begin{document}
\section*{Kourovka Notebook Problem 21.96}\label{2605}
{\small UnsolvedMath ID 2605 · KOU-21.96 · Group Theory · Kourovka Notebook - New Problems, Issue 21}
\subsection{Definitions and mathematical statement}
The order of an element $g$ of a group $G$ is the least positive integer $m$ with $g^m=1$, when one exists. A group is periodic, or torsion, if every element has finite order. It is locally finite if every subgroup generated by finitely many elements is finite; equivalently, each finite subset lies in a finite subgroup.

An involution is an element of order exactly two. For $x\in G$, its centralizer is the subgroup $C_G(x)=\{g\in G:gx=xg\}$. Suppose $G$ is periodic, contains an involution, and satisfies
\[
 \forall x\in G\text{ of even order},\qquad
 C_G(x)\text{ is locally finite}.
\]
Must $G$ itself be locally finite?

More explicitly, the centralizer hypothesis says that whenever $x$ has even order and $a_1,\ldots,a_r$ all commute with $x$, the subgroup $\langle a_1,\ldots,a_r\rangle$ is finite. The desired conclusion removes the requirement that these generators have a common even-order element in their centralizers. It asks for finiteness of every finitely generated subgroup, not for finiteness of $G$ as a whole. No common upper bound on element orders or on the sizes of these finite subgroups is assumed. Centralizers are tested for all even-order elements, not only involutions.
\subsection{Short English statement}
Must a periodic group with an involution be locally finite when every even-order element has a locally finite centralizer?
\subsection{Sources}
\begin{itemize}
\item[S1] ULAM.AI and the UnsolvedMath Contributors, supplied problems.json, record 2605 (KOU-21.96), Kourovka Notebook - New Problems, Issue 21. Catalogue identity and source transcription; CC BY 4.0.
\url{https://huggingface.co/datasets/ulamai/UnsolvedMath}.
\item[S2] The Kourovka Notebook, 21st edition, version 47 (30 September 2026), new problems, Problem 21.96. Expanded formulation of the supplied catalogue statement.
\url{https://arxiv.org/pdf/1401.0300v47}.
\end{itemize}
\end{document}
